\section{法律、宗教与继承：秩序的不同语言}
\label{sec:tang-imperial-order}

贞观元年修定律令、永徽四年完成《唐律疏议》，使初唐官府拥有较稳定的刑名与解释框架。法典的连续性不能被误读为社会的整齐：律文说明规范、疏议说明官方解释，而州县诉讼、身份登记和实际裁断仍须由文书与个案来观察。唐朝对隋制的继承，最可见于这种把既有制度重新编排为可援引语言的工作。%
\footnote{贞观律与《唐律疏议》见 ST-627-ZHENGUAN-CODE、ST-653-TANGLU-SHUYI，依据《旧唐书·刑法志》《唐律疏议》及《通典·刑法》。颁行与地方执行不可混同。}

635年阿罗本到长安获准译经，664年玄奘卒于玉华宫而译场与弟子网络仍在延续。这两件事不宜被写成唐廷对所有宗教一概“开放”的证明：许可、译经、寺院、施主与官府之间各有时段和区域的条件。它们只表明都城及其交通网络能够容纳多种文本、僧侣与仪式进入政治可见的空间；景教碑和佛教传记又各有自我说明与典范化的目的。%
\footnote{阿罗本与玄奘见 ST-635-NESTORIAN、ST-664-XUANZANG-DEATH，依据《旧唐书》相关记载、景教碑及《大慈恩寺三藏法师传》《续高僧传》等。入华、获准与具体社群规模须分开讨论。}

643年李承乾被废、649年李治即位，则提醒人们法令与宗教并不能消除皇室继承的风险。传记把谋反案写得层次分明，却难以恢复宫廷各派全部打算；较可靠的是太宗通过废立重新指定继承人，高宗即位后又须依托既有官僚和法制语言取得承认。帝国秩序不是单一制度的结果，而是在法庭、寺院、宫门和人事任命中不断被重申。%
\footnote{承乾案与高宗即位见 ST-643-CHENGQIAN、ST-649-TAIZONG-DEATH，依据《旧唐书·太宗本纪》《旧唐书·承乾传》《新唐书》相关纪传及《资治通鉴》。案情与人物动机受继承结果影响，不能作心理断言。}
